\scanpage{122}
\mspage{\handwritten{3}}
% Three pen strokes cross the reused typescript diagonally.
\deleted{$M_v(S')={}$ ensemble des $g\handwritten{\in}G(S')$ tels que \deleted{$g.W_{S'}$} % overtyped with xxxxx, xxxx
$v_{S'}\handwritten{\in}g.W(S')$, i.e. tels que $g^{-1}\,v_{S'}\handwritten{\in}W(S')$. Si $v$ est une section de $W$ et non seulement \handwritten{de} $V$, alors $M_v$ contient évidemment $N$.

Dans ces explicitations, on n'a pas utilisé le fait que $G,V,W$ étaient représentables (ni que le site sur lequel on travaille est \deleted{\illegible{1}} défini en termes de préschémas ! ). Mais supposons maintenant que $N$ soit représentable et fidèlement plat et quasi-compact sur $S$, et que $G/N$ soit représentable. Lorsque $S$ est le spectre d'un corps, et que $G$ est de type fini sur $k$, on sait que cette hypothèse est \deleted{\illegible{1}} nécessairement satisfaite (Exp VIII 6 \quad et Exp VI). On voit alors sur le diagramme cartésien $(\mathrm{xx})$, utilisant la théorie de la descente fidèlement plate quasi-compacte et le fait que \inserted{$Z\handwritten{\to}G\times_S V$} est une immersion fermée, que $X$ est représentable (il est obtenu par descente \deleted{\illegible{1}} du sous-préschéma fermé $Z$ de $G\times_S V$ par le morphisme fidèlement plat et quasi-compact $G\times_S V\handwritten{\to}G/N\times_S V$). Donc le diagramme (D) est un diagramme de morphismes de \deleted{\unsure{schémas}} prés\deleted{\unsure{c}}\handwritten{ch}émas sur $S$.

\margin{\handwritten{?}}Nous supposons par la suite que $S$ est le spectre d'un corps $k$, et que $G,V,W$ sont de type fini sur $k$. Soit \deleted{\illegible{1}} $\handwritten{\underline{\unsure{\mathfrak n|n}}}$ l\deleted{(}\handwritten{'}algèbre de Lie de $N$, donc on a $\dim N\handwritten{\leq}\operatorname{rang}\handwritten{\underline{\mathfrak n}}$, l'égalité étant vraie si et seulement si $N$ est lisse sur $k$ (Exp VI). \margin{\handwritten{?}}Soit $a\handwritten{\in}W(k)$, et considérons le sous-schéma $M_a$ de $G$ défini plus haut, contenant $N$, et isomorphe à $\unsure{\phi|\varnothing}^{-1}(a)$; nous désignerons par $\handwritten{\underline{\unsure{\mathfrak m|m}}}_a$ son espace \deleted{\unsure{d}}tangent de Zariski en l'élément neutre $e$ de $N$, de sorte qu'on a
\[
\handwritten{\underline{\mathfrak n}\subset\mathfrak m}_{a}\handwritten{\,,\qquad\deleted{\unsure{\dim}}\quad\dim N\leq\deleted{\unsure{\mathrm d}}\operatorname{rang}\mathfrak n\leq\operatorname{rang}\mathfrak m_a.}\tag*{\handwritten{(1)}}
\]

\uline{Lemme} 1.1. \deleted{\unsure{a)}} Avec les notations précédentes:

a) Considérons les conditions suivantes:

(i) $\handwritten{\mathfrak n}=\handwritten{\mathfrak m}_a$ \deleted{\illegible{1}} et $N$ est lisse sur $k$.}

% archive sheet number: 121
